\documentclass[12pt]{article}

% ---- Core math ----
\usepackage{amsmath}
\usepackage{amssymb}
\usepackage{amsthm}
\usepackage{mathtools}

% ---- Page layout ----
\usepackage[margin=1in]{geometry}

% ---- Graphics / diagrams ----
\usepackage{graphicx}
\usepackage{xcolor}
\usepackage{tikz}
\usetikzlibrary{arrows.meta,positioning,calc}
\usepackage{tikz-cd}
% The GrokRxiv sidebar uses the native LaTeX kernel shipout hook
% (\AddToHook{shipout/foreground}) rather than the legacy everypage package,
% avoiding deprecation warnings on current TeX distributions.

% ---- Tables ----
\usepackage{longtable}
\usepackage{booktabs}
\usepackage{array}
\usepackage{ragged2e}
% ragged-right text column: avoids underfull-hbox warnings in narrow p{} columns
\newcolumntype{L}[1]{>{\RaggedRight\arraybackslash}p{#1}}

% ---- References (hyperref before cleveref) ----
\usepackage{hyperref}
\hypersetup{
  colorlinks=true,
  linkcolor=blue!55!black,
  citecolor=green!45!black,
  urlcolor=blue!60!black
}
\usepackage{cleveref}

% ---- Theorem environments ----
\newtheorem{theorem}{Theorem}[section]
\newtheorem{proposition}[theorem]{Proposition}
\newtheorem{lemma}[theorem]{Lemma}
\newtheorem{corollary}[theorem]{Corollary}
\theoremstyle{definition}
\newtheorem{definition}[theorem]{Definition}
\newtheorem{example}[theorem]{Example}
\theoremstyle{remark}
\newtheorem{remark}[theorem]{Remark}

% ---- Notation macros ----
\newcommand{\Dom}{\mathsf{Dom}}
\newcommand{\Phys}{\mathsf{Phys}}
\newcommand{\Math}{\mathsf{Math}}
\newcommand{\Cat}{\mathsf{Cat}}
\newcommand{\Set}{\mathsf{Set}}
\newcommand{\Grpd}{\mathsf{Grpd}}
\newcommand{\Vect}{\mathsf{Vect}}
\newcommand{\FdHilb}{\mathsf{FdHilb}}
\newcommand{\Bord}{\mathsf{Bord}}
\newcommand{\Type}{\mathsf{Type}}
\newcommand{\Rep}{\mathsf{Rep}}
\newcommand{\Real}{\mathrm{Real}}
\newcommand{\Obs}{\mathrm{Obs}}
\newcommand{\Ob}{\operatorname{Ob}}
\newcommand{\Hom}{\operatorname{Hom}}
\newcommand{\id}{\operatorname{id}}
\newcommand{\Aut}{\operatorname{Aut}}
\newcommand{\Stab}{\operatorname{Stab}}
\newcommand{\Univ}{\mathcal{U}}
\newcommand{\Refl}{\mathsf{refl}}
\newcommand{\ua}{\mathsf{ua}}
\newcommand{\idtoequiv}{\mathsf{idtoequiv}}
\newcommand{\transport}{\mathsf{transport}}
\newcommand{\ap}{\mathsf{ap}}
\newcommand{\op}{\mathrm{op}}
\newcommand{\yon}{\mathsf{y}}
\DeclareMathOperator{\colim}{colim}
\newcommand{\llbracket}{\mathopen{[\![}}
\newcommand{\rrbracket}{\mathclose{]\!]}}
\newcommand{\sem}[1]{\llbracket #1 \rrbracket}

% ---- GrokRxiv DOI sidebar ----
\definecolor{grokgray}{RGB}{110,110,110}
\AddToHook{shipout/foreground}{%
  \ifnum\value{page}=1
    \begin{tikzpicture}[remember picture, overlay]
      \node[
        rotate=90,
        anchor=south,
        font=\Large\sffamily\bfseries\color{grokgray},
        inner sep=0pt
      ] at ([xshift=38pt, yshift=0.52\paperheight]current page.south west)
      {GrokRxiv:2026.07.category-theory-hott-composition\quad
       [\,math.CT\,]\quad
       04 Jul 2026};
    \end{tikzpicture}
  \fi
}

\title{\textbf{Category Theory and Homotopy Type Theory:\\ Composition and Identity}\\[0.6em]
\large Part V of \emph{A Math$\to$Physics Representation Library}}

\author{Matthew Long \\
\textit{The YonedaAI Collaboration} \\
\textit{YonedaAI Research Collective} \\
Chicago, IL \\
\texttt{research@yonedaai.com} $\cdot$ \url{https://yonedaai.com}}

\date{4 July 2026}

\begin{document}
\maketitle

\begin{abstract}
This is Part~V of a seven-part modular series that formalizes each mathematical
domain as a physical-representation module. Parts~I--IV built a
\emph{representation stack}: mathematical structures $M$ are sent by an
abstraction map $\Phi$ to abstract physical-information objects, realized by a
channel $\Real_\alpha$ into physical representations, and read off by an
observation map $\Obs$, giving the master pipeline
$\Obs_\alpha(M)=\Obs(\Real_\alpha(\Phi(M)))$. Those modules \emph{used}
functoriality (topological quantum field theory as a functor
$\Bord_n\to\Vect$; realization channels as translations between worlds) without
isolating the grammar that makes such translation coherent. The present module
supplies that grammar. We treat category theory as the theory of
\emph{composition and identity}: the associativity and unit laws that let
physical processes be chained, and functors and natural transformations as the
structure-preserving translations and their uniform deformations. We then treat
homotopy type theory (HoTT) as the theory of \emph{identity proper}: identity
types as paths, path induction as the sole generator of transport, and
Voevodsky's univalence axiom as the principle that equivalent representations
\emph{are} identified rather than merely related. Four theorems anchor the
module. (T1) Univalence turns Part~I's Axiom~(Equivalence) from a postulate into
a theorem for every internally definable realization. (T2) The
Curry--Howard--Lambek correspondence exhibits a single composite functor
$\mathsf{TypeTheory}\to\mathcal{C}\to\Phys$ interpreting logic, code, and
physical representation at once. (T3) In a dagger-compact category the existence
of uniform natural copying forces cartesianness, giving a structural no-cloning
theorem that explains why Part~VI's quantum high-availability encodings must be
genuine embeddings rather than replications. (T4) The groupoid (higher) quotient
$X/\!/G$ is strictly finer than the set quotient $X/G$, reconstructing the
``stacks retain gauge information'' phenomenon inside the truncation hierarchy of
HoTT and setting up Part~VI's stacks. Throughout we preserve the S/H/P epistemic
status discipline of the series, provide runnable Haskell encodings of the
category/functor/composition structures, and sketch an idiomatic Lean
formalization. Category theory and HoTT are thus the compositional glue of the
whole library: they are the module in which the laws of composition and identity
that bind Parts~I--VI are made explicit.
\end{abstract}

\tableofcontents

% =====================================================================
\section{Introduction}
\label{sec:intro}

\subsection{Where this module sits}

The series \emph{A Math$\to$Physics Representation Library} formalizes the thesis
that many physical quantities are not merely \emph{described by} mathematics but
are \emph{obtained as realizations} of structured mathematical information. The
organizing device is a \emph{representation stack} together with a
\emph{realization pipeline}
\begin{equation}
\label{eq:pipeline}
\begin{aligned}
M\in\Math_d \;\xrightarrow{\;\Phi\;}\; \text{abstract information }\Phi(M)
\;&\xrightarrow{\;\Real_\alpha\;}\; \text{physical representation}\\
&\xrightarrow{\;\Obs\;}\; \text{measurement data},
\end{aligned}
\end{equation}
compressed into the universal formula
\begin{equation}
\label{eq:master}
\boxed{\;\Obs_\alpha(M)=\Obs\bigl(\Real_\alpha(\Phi(M))\bigr)\;}
\quad\text{i.e.}\quad
\text{observable}=\Real(\text{abstract structure}).
\end{equation}
Each entry of the library carries an epistemic status label
$\sigma\in\{\mathrm{S},\mathrm{H},\mathrm{P}\}$: \textbf{S}~(standard
mathematics or mathematical physics), \textbf{H}~(strong heuristic dictionary
entry), \textbf{P}~(speculative ontological extension); composite labels
$\mathrm{S/H}$ and $\mathrm{H/P}$ denote entries that are standard as
mathematics but heuristic or speculative in their physical reading. This
labelling is part of the formalism, not decoration, and we apply it verbatim
throughout.

The modular ladder is
\[
\resizebox{\textwidth}{!}{$\displaystyle
\underbrace{\text{I}}_{\text{stack}}\to
\underbrace{\text{II}}_{\text{motives}}\to
\underbrace{\text{III}}_{\text{geometry}}\to
\underbrace{\text{IV}}_{\text{topology}}\to
\underbrace{\text{V}}_{\text{\bf category theory \& HoTT}}\to
\underbrace{\text{VI}}_{\text{stacks, gauge, quantum HA}}\to
\text{Synthesis}.$}
\]
Part~IV exhibited \emph{one} rigorous instance of ``physics as a functor'': a
topological quantum field theory (TQFT) is a symmetric monoidal functor
$Z:\Bord_n\to\mathcal{C}$ \cite{Atiyah1988,Lurie2009}. But Part~IV used the word
``functor'' without developing the general theory of functors, and it used
``equivalent realizations describe the same physics'' without a formal notion of
identity of structures. Part~V is exactly the abstraction step that provides
both:
\begin{center}
\emph{category theory is the grammar of physical composition;\\
HoTT is the logic of spaces, paths, equivalences, and higher identity.}
\end{center}
Because every realization channel $\Real_\alpha$, every duality, and every
symmetry of Parts~I--IV is a morphism, functor, or equivalence, Part~V is the
\emph{compositional glue} of the whole library: it is where the laws of
composition and identity that bind Parts~I--VI are isolated and proved. It also
sets up Part~VI, whose stacks are groupoid-valued sheaves built from the
categorical and truncation machinery introduced here.

\subsection{Contributions}

\begin{enumerate}
\item \textbf{A categorical recasting of the realization pipeline}
(\cref{sec:framework}). We show that \eqref{eq:pipeline} is literally a composite
of functors between the groupoids of a representation prestack, so that the
associativity and unit laws of the library are the associativity and unit laws
of a category (\Cref{thm:pipeline-functor}). This turns Part~I's informal
``composition of compatible translations'' into a theorem.

\item \textbf{The composition/identity core}
(\cref{sec:cat,sec:monoidal}). We give precise definitions of categories,
functors, natural transformations, functor categories, and monoidal and
dagger-compact structure, with the interchange law and an Eckmann--Hilton
argument (\Cref{prop:eckmann-hilton}) that already exhibits the interaction of
two compositions with a shared identity---the thematic centre of the module.

\item \textbf{Univalence discharges the Equivalence axiom}
(\Cref{thm:univalence}, T1). We prove that any internally definable realization
respects equivalence of structures, so Part~I's Axiom~(Equivalence) becomes a
theorem in a univalent model.

\item \textbf{A single functor for logic, code, and physics}
(\Cref{thm:chl}, T2). The Curry--Howard--Lambek correspondence yields a composite
functor $\mathsf{TypeTheory}\to\mathcal{C}\to\Phys$ interpreting the three-way
dictionary (proof $=$ program $=$ physical realization) as one categorical
translation.

\item \textbf{A structural no-cloning theorem} (\Cref{thm:nocloning}, T3). In a
dagger-compact category, uniform natural copying forces the monoidal product to
be a categorical product; since the tensor product of finite-dimensional Hilbert
spaces is not a product, quantum high availability (Part~VI) is
\emph{representationally forced} to be encoding, not replication.

\item \textbf{Groupoid quotients are finer than set quotients}
(\Cref{thm:truncation}, T4). We reconstruct ``stacks retain gauge information''
purely inside the HoTT truncation hierarchy: $X/\!/G$ is a $1$-type retaining
stabilizer loops, while $X/G$ is its $0$-truncation, forgetting them.

\item \textbf{Formal artifacts.} A runnable Haskell package encoding categories,
functors, natural transformations, dagger-compact structure, and
groupoid-vs-set quotients (\cref{sec:haskell}), and an idiomatic Lean sketch of
the category/functor/identity-law types (\cref{sec:lean}).
\end{enumerate}

\subsection{Epistemic status and limitations}

We reproduce the series-wide limitations in the form relevant to this module.
Category theory and HoTT are \emph{status~S mathematics}; their status becomes
$\mathrm{S/H}$ or $\mathrm{H}$ only when we read a specific categorical or
type-theoretic construction as a claim \emph{about physics}. In particular: (i)
we do \emph{not} claim a single universal functor $\Math\to\Phys$---the
representation approach is deliberately domain-indexed
(\cref{sec:framework}); (ii) reading a type $A$ as a ``space of physical
possibilities'' is a heuristic semantic layer (status~$\mathrm{H}$), flagged as
such; (iii) univalence is a theorem-generating principle \emph{about
mathematical structures}, and its physical reading (``equivalent representations
are the same physical content'') is $\mathrm{S/H}$. These flags are carried
explicitly in the dictionary of \cref{sec:dictionary} and \cref{app:library}.

\subsection{Outline}

\Cref{sec:framework} recalls the representation-stack framework and casts the
pipeline categorically. \Cref{sec:cat} develops categories, functors, natural
transformations, and the interchange/Eckmann--Hilton phenomena.
\Cref{sec:monoidal} develops monoidal, symmetric, compact-closed and
dagger-compact structure and string diagrams. \Cref{sec:hott} develops HoTT:
identity types, path induction, transport, equivalence, and univalence.
\Cref{sec:results} proves the four theorems T1--T4. \Cref{sec:dictionary}
presents the category-theory/HoTT representation dictionary with S/H/P labels.
\Cref{sec:haskell,sec:lean} give the Haskell and Lean formalizations.
\Cref{sec:discussion} discusses composition in the ladder and limitations;
\cref{sec:conclusion} concludes. \Cref{app:library} contains the full
dictionary longtables.

% =====================================================================
\section{Mathematical framework: the representation stack, categorically}
\label{sec:framework}

We recall the minimum from Part~I and immediately reorganize it so that the
categorical language of \cref{sec:cat} applies.

\begin{definition}[Representation entry, {Part~I}]
\label{def:entry}
A \emph{representation entry} is a quadruple $E=(M,P,\tau,\sigma)$ where
$M\in\Math_d$ is a mathematical structure in a domain $d\in\Dom$, $P\in\Phys$ is
a proposed physical representation, $\tau$ is a semantic translation datum (a
functor, a natural transformation, an equivalence class of models, or a weaker
interpretation map), and $\sigma\in\{\mathrm{S},\mathrm{H},\mathrm{P}\}$ is the
epistemic status label.
\end{definition}

\begin{definition}[Representation prestack, {Part~I}]
\label{def:prestack}
A \emph{representation prestack} is a pseudofunctor
$\Rep_{\Phys}:\Dom^{\op}\to\Grpd$ assigning to each domain $d$ a groupoid
$\Rep_{\Phys}(d)$ of representation entries, and to a domain refinement
$u:e\to d$ a restriction functor $u^{*}:\Rep_{\Phys}(d)\to\Rep_{\Phys}(e)$,
functorial up to coherent natural isomorphism.
\end{definition}

The words ``groupoid'', ``functor'', ``natural isomorphism'', and
``pseudofunctor'' in \cref{def:prestack} are precisely the notions
\cref{sec:cat} makes rigorous; Part~I used them on trust. We keep the realization
pipeline \eqref{eq:pipeline} but now name its stages as morphisms in categories.

\begin{definition}[Realization data]
\label{def:realization}
Fix a domain $d$. A \emph{realization datum} of level $\alpha$ consists of:
an \emph{abstraction} functor $\Phi:\Math_d\to\mathcal{A}$ into a category
$\mathcal{A}$ of abstract physical-information objects; a \emph{realization}
functor $\Real_\alpha:\mathcal{A}\to\Phys$; and an \emph{observation} functor
$\Obs:\Phys\to\mathcal{M}$ into a category $\mathcal{M}$ of measurement data. The
\emph{pipeline of level $\alpha$} is the composite functor
\[
\Obs_\alpha \;:=\; \Obs\circ\Real_\alpha\circ\Phi \;:\; \Math_d\longrightarrow\mathcal{M}.
\]
\end{definition}

\begin{theorem}[The pipeline is a composite of functors]
\label{thm:pipeline-functor}
With the data of \cref{def:realization}, $\Obs_\alpha$ is a functor, and its
formation is associative and unital: for composable domain refinements the
restriction functors compose (up to coherent isomorphism), and the identity
refinement acts as the identity functor. Consequently the master formula
\eqref{eq:master} is the object-level shadow of a genuine composition of
morphisms in $\Cat$, and Part~I's closure operation ``composition of compatible
translations'' is the composition law of the category $\Cat$ of categories.
\end{theorem}

\begin{proof}
Functors compose: if $F:\mathcal{X}\to\mathcal{Y}$ and $G:\mathcal{Y}\to\mathcal{Z}$
preserve identities and composition, then $(G\circ F)(\id_X)=G(\id_{FX})=\id_{GFX}$
and $(G\circ F)(g\circ f)=G(Fg\circ Ff)=G(Fg)\circ G(Ff)$, so $G\circ F$ is a
functor. Applying this twice to $\Phi$, $\Real_\alpha$, $\Obs$ gives that
$\Obs_\alpha$ is a functor. Associativity of $\circ$ in $\Cat$ is inherited from
associativity of function composition on objects and on hom-sets; the identity
functor $\id_{\mathcal{X}}$ is a two-sided unit for the same reason. The
pseudofunctoriality clause of \cref{def:prestack} states exactly that
$u^{*}\circ v^{*}\cong (v\circ u)^{*}$ and $(\id_d)^{*}\cong\id$, which is
associativity and unitality up to coherent isomorphism. Reading
\eqref{eq:master} on objects, $\Obs_\alpha(M)=\Obs(\Real_\alpha(\Phi(M)))$ is the
value of the composite functor at $M$, i.e.\ the object component of the
composition $\Obs\circ\Real_\alpha\circ\Phi$ in $\Cat$.
\end{proof}

\Cref{thm:pipeline-functor} is the reason category theory belongs at Level~V of
the ladder: the entire library was already, implicitly, a diagram in $\Cat$. The
remaining sections make the ingredients precise and extract the four physical
consequences T1--T4. The pipeline as a commuting diagram is:
\begin{equation}
\label{eq:pipeline-cd}
\begin{tikzcd}[column sep=large, row sep=large]
\Math_d \arrow[r, "\Phi"] \arrow[rrr, bend left=22, "\Obs_\alpha"]
& \mathcal{A} \arrow[r, "\Real_\alpha"]
& \Phys \arrow[r, "\Obs"]
& \mathcal{M} \\
\end{tikzcd}
\end{equation}

% =====================================================================
\section{Categories, functors, and natural transformations}
\label{sec:cat}

\subsection{Categories: composition and identity}

\begin{definition}[Category]
\label{def:category}
A \emph{category} $\mathcal{C}$ consists of a collection $\Ob(\mathcal{C})$ of
\emph{objects}; for each ordered pair $(A,B)$ a set $\Hom_{\mathcal{C}}(A,B)$ of
\emph{morphisms}; a \emph{composition} operation
$\circ:\Hom(B,C)\times\Hom(A,B)\to\Hom(A,C)$; and for each object $A$ an
\emph{identity} morphism $\id_A\in\Hom(A,A)$, subject to the two axioms
\begin{align}
\text{(associativity)}\qquad & h\circ(g\circ f)=(h\circ g)\circ f, \label{eq:assoc}\\
\text{(unit)}\qquad & \id_B\circ f=f=f\circ\id_A \quad\text{for } f:A\to B.
\label{eq:unit}
\end{align}
\end{definition}

Physically (status $\mathrm{S/H}$), objects are systems or state spaces,
morphisms $f:A\to B$ are processes or transitions, $g\circ f$ is ``do $f$, then
$g$'', and $\id_A$ is the do-nothing process. \eqref{eq:assoc} says that grouping
of a process chain is physically immaterial; \eqref{eq:unit} says that inserting
a trivial step changes nothing. These two laws are the entire content of
``composition and identity'' at the $1$-categorical level.

\begin{example}[Groups and representations as one-object categories]
\label{ex:BG}
A monoid $M$ is a category with one object $\ast$ and $\Hom(\ast,\ast)=M$, with
composition the monoid product and identity the unit. A group $G$ is such a
category in which every morphism is invertible---a one-object groupoid, written
$BG$ (the \emph{delooping}). A linear representation $\rho:G\to GL(V)$ is exactly
a functor $BG\to\Vect$: it sends the unique object to $V$ and each group element
to an invertible linear map, and functoriality $\rho(gh)=\rho(g)\rho(h)$,
$\rho(e)=\id_V$ is the definition of a representation. This is the first place
the slogan ``a physical theory is a functor'' becomes literal: a symmetry action
\emph{is} a functor out of a delooping.
\end{example}

\begin{example}[Preorders, spaces, and $\FdHilb$]
\label{ex:examples}
(i) A preorder $(P,\le)$ is a category with at most one morphism $A\to B$,
present iff $A\le B$; composition is transitivity, identities are reflexivity.
(ii) $\Set$ has sets as objects and functions as morphisms; $\Vect$ has vector
spaces and linear maps; $\FdHilb$ has finite-dimensional Hilbert spaces and
linear maps. (iii) The bordism category $\Bord_n$ has closed $(n{-}1)$-manifolds
as objects and $n$-dimensional cobordisms (up to diffeomorphism) as morphisms; a
TQFT is a functor $\Bord_n\to\Vect$ (\cref{ex:examples} feeds
\cref{sec:monoidal}).
\end{example}

\begin{definition}[Isomorphism, groupoid]
A morphism $f:A\to B$ is an \emph{isomorphism} if there is $g:B\to A$ with
$g\circ f=\id_A$ and $f\circ g=\id_B$; then $g$ is unique and written $f^{-1}$. A
\emph{groupoid} is a category in which every morphism is an isomorphism.
Physically (status $\mathrm{S}$) an isomorphism is a reversible equivalence:
``the same structure in different coordinates.''
\end{definition}

\subsection{Functors: translation between worlds}

\begin{definition}[Functor]
\label{def:functor}
A \emph{functor} $F:\mathcal{C}\to\mathcal{D}$ assigns to each object $A$ an
object $FA$ and to each morphism $f:A\to B$ a morphism $Ff:FA\to FB$, preserving
identities and composition:
\begin{equation}
\label{eq:functor-laws}
F(\id_A)=\id_{FA},\qquad F(g\circ f)=Fg\circ Ff.
\end{equation}
\end{definition}

A functor is the structure-preserving translation of the composition/identity
laws from one world to another; a physical theory, a quantization, or a
realization channel $\Real_\alpha$ is such a translation (status~$\mathrm{S}$ for
TQFT; $\mathrm{S/H}$ for a general realization functor). \eqref{eq:functor-laws}
is precisely what makes \eqref{eq:pipeline-cd} commute stagewise.

\begin{definition}[Natural transformation]
\label{def:nat}
Given functors $F,G:\mathcal{C}\to\mathcal{D}$, a \emph{natural transformation}
$\eta:F\Rightarrow G$ is a family of morphisms $\eta_A:FA\to GA$ (its
\emph{components}) such that for every $f:A\to B$ the \emph{naturality square}
commutes:
\begin{equation}
\label{eq:naturality}
\begin{tikzcd}[column sep=large, row sep=large]
FA \arrow[r, "\eta_A"] \arrow[d, "Ff"'] & GA \arrow[d, "Gf"] \\
FB \arrow[r, "\eta_B"'] & GB
\end{tikzcd}
\qquad\Longleftrightarrow\qquad
Gf\circ\eta_A=\eta_B\circ Ff.
\end{equation}
If each $\eta_A$ is an isomorphism, $\eta$ is a \emph{natural isomorphism} and
$F\cong G$. Physically (status~$\mathrm{S/H}$) a natural transformation is a
uniform deformation of one translation into another---a field redefinition or
theory equivalence that acts coherently on all systems at once.
\end{definition}

\begin{definition}[Functor category]
\label{def:functor-cat}
For categories $\mathcal{C},\mathcal{D}$, the \emph{functor category}
$[\mathcal{C},\mathcal{D}]$ has functors $\mathcal{C}\to\mathcal{D}$ as objects
and natural transformations as morphisms, with \emph{vertical composition}
$(\theta\cdot\eta)_A=\theta_A\circ\eta_A$ and identity the identity natural
transformation. There is also \emph{horizontal composition} $\ast$ of natural
transformations along a composite of functors.
\end{definition}

\subsection{Interchange and Eckmann--Hilton}

In the functor $2$-category $\Cat$, objects are categories, $1$-morphisms are
functors, and $2$-morphisms are natural transformations. These $2$-morphisms
carry two composition laws. Their compatibility is the single most important
coherence in the module, because ``composition and identity'' now appears
twice---two compositions sharing one family of identities.

\begin{proposition}[Interchange law]
\label{prop:interchange}
In any $2$-category, for $2$-morphisms arranged so that both sides are defined,
\begin{equation}
\label{eq:interchange}
(\theta'\cdot\theta)\ast(\eta'\cdot\eta)=(\theta'\ast\eta')\cdot(\theta\ast\eta),
\end{equation}
where $\cdot$ is vertical and $\ast$ is horizontal composition. In particular
horizontal and vertical composition of natural transformations satisfy
\eqref{eq:interchange}.
\end{proposition}

\begin{proof}
Both sides are computed by evaluating components and using naturality
\eqref{eq:naturality}. Writing the horizontal composite of $\eta:F\Rightarrow G$
(along $\mathcal{C}\to\mathcal{D}$) with $\eta':F'\Rightarrow G'$ (along
$\mathcal{D}\to\mathcal{E}$) as $(\eta'\ast\eta)_A=\eta'_{GA}\circ F'\eta_A
=G'\eta_A\circ\eta'_{FA}$ (the two expressions agree by naturality of $\eta'$ at
$\eta_A$), one expands both sides of \eqref{eq:interchange} on an object $A$ and
finds each equal to the common diagonal of the evident commuting square; the
equality is forced by the naturality of the four transformations. A full
component computation is standard \cite[Ch.~II]{MacLane1998}.
\end{proof}

\begin{proposition}[Eckmann--Hilton]
\label{prop:eckmann-hilton}
Let a set $S$ carry two unital binary operations $\cdot$ and $\ast$ sharing a
common unit $e$ and satisfying the interchange law
$(a\ast b)\cdot(c\ast d)=(a\cdot c)\ast(b\cdot d)$. Then the two operations
coincide, are associative, and are commutative.
\end{proposition}

\begin{proof}
First $e=e\ast e=e\cdot e$. Interchange gives, for all $a,b$,
\[
a\ast b=(a\cdot e)\ast(e\cdot b)=(a\ast e)\cdot(e\ast b)=a\cdot b,
\]
so $\ast=\cdot$; call it $\times$. Then
$a\times b=(e\ast a)\times(b\ast e)=(e\times b)\ast(a\times e)=b\times a$, giving
commutativity, and interchange collapses to associativity.
\end{proof}

\begin{remark}[Why this belongs to a module titled ``composition and identity'']
\label{rem:eh}
\Cref{prop:eckmann-hilton} is the algebraic reason that the second homotopy
group of a space is abelian, and in HoTT it is the statement that the higher
identity types $\Omega^2$ (paths between paths) are commutative: the two ways to
compose $2$-dimensional identities---horizontally and vertically---share the
constant identity and interchange, hence agree and commute. Thus ``composition''
and ``identity'' are not two topics but one: at each new dimension the identity
laws constrain how compositions may interact, and Eckmann--Hilton is the first
nontrivial constraint. This directly forecasts the higher-identity content of
\cref{sec:hott}.
\end{remark}

\subsection{Yoneda: a system is known by its probes}

\begin{lemma}[Yoneda]
\label{lem:yoneda}
For a locally small category $\mathcal{C}$, an object $A\in\mathcal{C}$, and a
functor $F:\mathcal{C}^{\op}\to\Set$, there is a bijection natural in both $A$
and $F$,
\begin{equation}
\label{eq:yoneda}
\mathrm{Nat}\bigl(\Hom_{\mathcal{C}}(-,A),\,F\bigr)\;\cong\;F(A),
\qquad \eta\mapsto\eta_A(\id_A).
\end{equation}
In particular the \emph{Yoneda embedding} $\yon\!:\mathcal{C}\hookrightarrow
[\mathcal{C}^{\op},\Set]$, $A\mapsto\Hom_{\mathcal{C}}(-,A)$, is full and
faithful, so $A\cong B$ iff $\Hom(-,A)\cong\Hom(-,B)$.
\end{lemma}

\begin{proof}
Given a natural transformation $\eta:\Hom(-,A)\Rightarrow F$, set
$u:=\eta_A(\id_A)\in F(A)$. Naturality at any $f:X\to A$ (viewed as an element of
$\Hom(X,A)$) forces $\eta_X(f)=F(f)(u)$, so $\eta$ is determined by $u$; and any
$u\in F(A)$ defines a natural $\eta$ by this formula. The two assignments are
mutually inverse, giving \eqref{eq:yoneda}. Naturality in $A,F$ is a direct
check. Taking $F=\Hom(-,B)$ yields
$\mathrm{Nat}(\Hom(-,A),\Hom(-,B))\cong\Hom(A,B)$, i.e.\ full faithfulness of
$\yon$; reflection of isomorphisms follows because full and faithful functors
reflect isomorphisms.
\end{proof}

Physically (status~$\mathrm{S/H}$) the Yoneda lemma is the statement that a
system is completely determined by the totality of its responses to all probes:
knowing $\Hom(-,A)$---all ways of mapping into $A$---determines $A$ up to
isomorphism. This is the categorical form of an operational/measurement-based
identity criterion and reappears as the ``univalence for representables'' motif
in \cref{sec:hott}.

% =====================================================================
\section{Monoidal and dagger-compact structure}
\label{sec:monoidal}

Parallel composition of systems---putting two experiments side by side---is not
captured by $\circ$ alone. It is captured by a tensor product, and the reversal
of processes by a dagger. These structures make $\FdHilb$ into the ambient
category of finite quantum mechanics and give the string-diagram calculus used
in Part~VI.

\begin{definition}[Monoidal category]
\label{def:monoidal}
A \emph{monoidal category} is a category $\mathcal{C}$ with a functor
$\otimes:\mathcal{C}\times\mathcal{C}\to\mathcal{C}$, a unit object $I$, and
natural isomorphisms $\alpha_{A,B,C}:(A\otimes B)\otimes C\to A\otimes(B\otimes
C)$, $\lambda_A:I\otimes A\to A$, $\rho_A:A\otimes I\to A$ satisfying Mac Lane's
pentagon and triangle coherence axioms. It is \emph{symmetric} if it also has a
natural isomorphism $\sigma_{A,B}:A\otimes B\to B\otimes A$ with
$\sigma_{B,A}\circ\sigma_{A,B}=\id$ and the hexagon coherence.
\end{definition}

That $\otimes$ is a \emph{bifunctor} means it respects composition in each
argument and satisfies the interchange law $(g\otimes g')\circ(f\otimes
f')=(g\circ f)\otimes(g'\circ f')$: sequential and parallel composition of
processes commute, which is the algebra of quantum circuits. Physically
(status~$\mathrm{S}$) $A\otimes B$ is the composite system, $I$ is the trivial
system, and $\sigma$ is the swap---bosonic-like exchange when it squares to the
identity, anyonic braiding when it does not.

\begin{definition}[Compact closed and dagger-compact category]
\label{def:dagger}
A \emph{compact closed} category is a symmetric monoidal category in which every
object $A$ has a dual $A^{*}$ equipped with a \emph{unit}
$\eta_A:I\to A^{*}\otimes A$ (coevaluation, a ``cup'') and \emph{counit}
$\varepsilon_A:A\otimes A^{*}\to I$ (evaluation, a ``cap'') satisfying the
\emph{snake} (zig-zag) equations
\begin{equation}
\label{eq:snake}
(\varepsilon_A\otimes\id_A)\circ(\id_A\otimes\eta_A)=\id_A,\qquad
(\id_{A^{*}}\otimes\varepsilon_A)\circ(\eta_A\otimes\id_{A^{*}})=\id_{A^{*}}.
\end{equation}
A \emph{dagger-compact} category adds an involutive, identity-on-objects,
contravariant functor $(-)^{\dagger}$ that is monoidal and compatible with the
compact structure, so that $\varepsilon_A=\eta_A^{\dagger}\circ\sigma$ up to the
canonical isomorphisms. Physically $(-)^{\dagger}$ is the adjoint/time-reversal.
\end{definition}

\begin{example}[$\FdHilb$ as the arena of finite quantum mechanics]
\label{ex:fdhilb}
$\FdHilb$ is dagger-compact: $\otimes$ is the tensor product of Hilbert spaces,
$I=\mathbb{C}$, $A^{*}$ is the dual space, $\eta_A(1)=\sum_i e_i\otimes e_i$ is
the (unnormalized) Bell state, $\varepsilon_A$ is the pairing, and
$(-)^{\dagger}$ is the Hermitian adjoint. The snake equations \eqref{eq:snake}
are the statement that a maximally entangled pair can be created and then
partially annihilated to yield an identity wire---the diagrammatic engine of
quantum teleportation \cite{AbramskyCoecke2004,BaezStay2009}.
\end{example}

The snake equation admits the string-diagram reading in which a cup followed by a
cap straightens a bent wire:
\begin{equation}
\label{eq:snake-cd}
\begin{tikzcd}[column sep=large, row sep=large]
A \arrow[r, "\id_A\otimes\eta_A"] \arrow[dr, "\id_A"'] & A\otimes A^{*}\otimes A \arrow[d, "\varepsilon_A\otimes\id_A"] \\
& A
\end{tikzcd}
\end{equation}

% =====================================================================
\section{Homotopy type theory: identity as path}
\label{sec:hott}

We now pass from composition (category theory) to identity proper. HoTT is a
dependent type theory in which the identity type is interpreted homotopically and
Voevodsky's univalence axiom governs identity of structures.

\subsection{Types, terms, and dependent types}

We use the judgement $a:A$ (``$a$ is a term of type $A$''). A \emph{dependent
type} $B:A\to\Univ$ (a type family) assigns a type $B(x)$ to each $x:A$. From it
we form the dependent pair and function types
\begin{equation}
\label{eq:sigmapi}
\sum_{x:A}B(x)\quad(\text{total space}),\qquad
\prod_{x:A}B(x)\quad(\text{space of sections}).
\end{equation}
Physically (status~$\mathrm{S/H}$), $A$ is a space of states or possibilities,
$a:A$ a witnessing state, $B(x)$ a fibre of field data over the base point $x$,
$\sum_{x:A}B(x)$ the total configuration space, and $\prod_{x:A}B(x)$ the space
of fields (sections) assigning data at every point.

\subsection{Identity types and path induction}

\begin{definition}[Identity type]
\label{def:idtype}
For $A:\Univ$ and $a,b:A$ there is a type $\mathsf{Id}_A(a,b)$, written $a=b$,
whose terms are \emph{identifications} (paths) from $a$ to $b$. It is generated
by the constructor $\Refl_a:a=a$, subject to the following induction principle.
\end{definition}

\begin{definition}[Path induction, the J-rule]
\label{def:J}
To define a section of a family $C:\prod_{x,y:A}(x=y)\to\Univ$ it suffices to
define it on reflexivity: given
$c:\prod_{x:A}C(x,x,\Refl_x)$ there is a unique (up to identification) term
\[
\mathsf{J}(c):\prod_{x,y:A}\prod_{p:x=y}C(x,y,p)
\qquad\text{with}\qquad \mathsf{J}(c)(x,x,\Refl_x)\equiv c(x).
\]
Path induction is the \emph{only} elimination rule for identity types.
\end{definition}

\begin{lemma}[Transport]
\label{lem:transport}
For any $B:A\to\Univ$ and $p:a=b$ there is an equivalence
$\transport^{B}(p):B(a)\xrightarrow{\ \sim\ }B(b)$, with
$\transport^{B}(\Refl_a)=\id_{B(a)}$, and $\transport^{B}$ is functorial:
$\transport^{B}(q\cdot p)=\transport^{B}(q)\circ\transport^{B}(p)$ for
$p:a=b$, $q:b=c$.
\end{lemma}

\begin{proof}
Apply \cref{def:J} to the family $C(x,y,p):=\bigl(B(x)\to B(y)\bigr)$ with
$c(x):=\id_{B(x)}$; this yields $\transport^{B}(p):B(a)\to B(b)$ with
$\transport^{B}(\Refl_a)\equiv\id$. That each $\transport^{B}(p)$ is an
equivalence follows by path induction on $p$ (it holds at $\Refl$, hence
everywhere), with inverse $\transport^{B}(p^{-1})$. Functoriality is again path
induction: fix $p$ and induct on $q$; the case $q=\Refl_b$ reduces to
$\transport^{B}(p)=\transport^{B}(p)\circ\id$, which holds definitionally.
\end{proof}

\begin{lemma}[Action on paths]
\label{lem:ap}
Every function $f:A\to B$ acts on identifications: there is
$\ap_f:(a=a')\to(f(a)=f(a'))$ with $\ap_f(\Refl_a)=\Refl_{f(a)}$, respecting
inverses and concatenation. Thus functions automatically preserve identity---the
type-theoretic counterpart of functoriality \eqref{eq:functor-laws}.
\end{lemma}

\begin{proof}
Path induction on the family $C(x,y,p):=(f(x)=f(y))$ with $c(x):=\Refl_{f(x)}$.
Preservation of concatenation and inverses is a further path induction.
\end{proof}

\subsection{Equivalence and univalence}

\begin{definition}[Equivalence]
\label{def:equiv}
A function $f:A\to B$ \emph{is an equivalence}, written $\mathsf{isEquiv}(f)$, if
its fibres are contractible; equivalently if it has a two-sided
homotopy-inverse. The type of equivalences is
$A\simeq B:=\sum_{f:A\to B}\mathsf{isEquiv}(f)$. The canonical map
\begin{equation}
\label{eq:idtoequiv}
\idtoequiv:(A=B)\longrightarrow(A\simeq B)
\end{equation}
sends $\Refl_A$ to the identity equivalence (by path induction).
\end{definition}

\begin{definition}[Univalence axiom, Voevodsky]
\label{def:univalence}
A universe $\Univ$ is \emph{univalent} if for all $A,B:\Univ$ the map
$\idtoequiv$ of \eqref{eq:idtoequiv} is itself an equivalence:
\begin{equation}
\label{eq:ua}
\boxed{\;(A=B)\;\simeq\;(A\simeq B)\;.}
\end{equation}
Its inverse is written $\ua:(A\simeq B)\to(A=B)$.
\end{definition}

Univalence is the exact formal content of the representation-theoretic slogan
``equivalent representations can be treated as the same physical content'':
identity of structures and equivalence of structures are themselves identified.
It is used in \cref{sec:results} to discharge Part~I's Axiom~(Equivalence).

\begin{definition}[Truncation levels]
\label{def:trunc}
A type $A$ is \emph{contractible} (level $-2$) if
$\sum_{a:A}\prod_{x:A}(a=x)$ is inhabited; a \emph{mere proposition}
(level $-1$) if all its terms are identified; a \emph{set}
(level $0$) if each identity type $a=b$ is a mere proposition; a
\emph{groupoid} (level $1$) if each $a=b$ is a set; and so on up the
$n$-truncation hierarchy. Each level $n$ carries an idempotent truncation
modality $\lVert-\rVert_n$ universally forcing a type to level $n$.
\end{definition}

\begin{example}[The circle]
\label{ex:circle}
The higher inductive type $S^1$ is generated by a point $\mathsf{base}:S^1$ and a
path $\mathsf{loop}:\mathsf{base}=\mathsf{base}$. It is a $1$-type, not a set:
its loop space is $\Omega S^1=(\mathsf{base}=\mathsf{base})\simeq\mathbb{Z}$, so
$\pi_1(S^1)=\mathbb{Z}$ is provable synthetically. Physically (status~$\mathrm{H}$)
$S^1$ models a phase circle or winding sector, and its integer loop space is the
winding number/topological charge.
\end{example}

% =====================================================================
\section{Results: the four theorems}
\label{sec:results}

We now prove the four theorems that give this module its physical payload. T1 and
T2 concern identity and translation; T3 and T4 concern quantum information and
gauge, feeding Part~VI.

\subsection{T1: univalence discharges the Equivalence axiom}

\begin{theorem}[Univalence respects realization]
\label{thm:univalence}
Let $\Univ$ be a univalent universe modelling Part~I's $\Math_d$, and let
$\Real:\Univ\to\Phys$ be any \emph{internally definable} realization, i.e.\ a
function definable in the type theory (equivalently, a term
$\Real:\prod_{X:\Univ}\Phys$). Then for all $A,B:\Univ$,
\[
(A\simeq B)\ \longrightarrow\ \bigl(\Real(A)=\Real(B)\bigr),
\]
and moreover $\transport^{\Real}$ along the induced path is an equivalence
$\Real(A)\simeq\Real(B)$. Hence any internally definable realization assigns the
same physical representation to equivalent mathematical structures: Part~I's
Axiom~(Equivalence) is a \emph{theorem}, not a postulate.
\end{theorem}

\begin{proof}
Let $e:A\simeq B$. By univalence \eqref{eq:ua}, $\ua(e):A=B$ is an
identification of $A$ and $B$ in $\Univ$. Apply \cref{lem:ap} to $\Real$:
$\ap_{\Real}(\ua(e)):\Real(A)=\Real(B)$, which is the required identification of
physical representations. To obtain an equivalence we transport along the path
$\ua(e):A=B$ \emph{in the base $\Univ$}, using the type family
$X\mapsto\Real(X)$ (with $\Phys$ read as a universe of types): by
\cref{lem:transport}, $\transport^{X\mapsto\Real(X)}(\ua(e)):\Real(A)\to
\Real(B)$ is an equivalence, invertible with inverse the transport along
$\ua(e)^{-1}=\ua(e^{-1})$. (Equivalently, applying $\idtoequiv$ to the path
$\ap_{\Real}(\ua(e))$ in $\Phys$ yields the same equivalence directly.) Since
$\ap_{\Real}(\Refl_A)=\Refl_{\Real(A)}$,
the construction sends the identity equivalence to reflexivity, so it is the
identity on the diagonal, as required for a bona fide discharge of the axiom.
\end{proof}

\begin{remark}[Gauge-dependent vs.\ gauge-invariant]
\label{rem:gauge-invariance}
\Cref{thm:univalence} applies only to \emph{internally definable} realizations.
A realization defined by an \emph{external} choice---picking a basis, a gauge, or
a coordinate patch---need not be a term of the type theory and may fail to
respect equivalence. This is exactly the physical distinction between
gauge-invariant quantities (internally definable; equivalence-respecting) and
gauge-dependent quantities (externally defined; not). Univalence thereby draws
the gauge-invariance line as a definability condition, anticipating Part~VI's
treatment of gauge redundancy.
\end{remark}

\subsection{T2: one functor for logic, code, and physics}

\begin{theorem}[Curry--Howard--Lambek semantics of the pipeline]
\label{thm:chl}
Let $T$ be a Martin-L\"of dependent type theory with $\Sigma$-, $\Pi$-, and
identity types, and let $\mathrm{Syn}(T)$ be its syntactic category (objects:
contexts; morphisms: substitutions up to definitional equality). For any locally
cartesian closed category $\mathcal{C}$ modelling $T$ there is an interpretation
functor
\[
\llbracket-\rrbracket:\mathrm{Syn}(T)\longrightarrow\mathcal{C}
\]
sending types to objects and terms to (global) sections, and it is sound: it
preserves the type-forming operations up to canonical isomorphism. Composing with
a realization functor $R:\mathcal{C}\to\Phys$ yields a single composite
\[
\mathrm{Syn}(T)\xrightarrow{\ \llbracket-\rrbracket\ }\mathcal{C}
\xrightarrow{\ R\ }\Phys
\]
that interprets, simultaneously and by one categorical translation, the logical
reading (propositions as types, proofs as terms), the computational reading
(programs as terms, evaluation as normalization), and the physical reading
(state spaces as types, processes as terms, realization as $R$).
\end{theorem}

\begin{proof}
The existence and soundness of $\llbracket-\rrbracket$ is the
Curry--Howard--Lambek/Seely correspondence between Martin-L\"of type theory with
$\Sigma,\Pi$, and identity types and locally cartesian closed categories
\cite{LambekScott1986,MacLane1998,HoTT2013}: $\Sigma$-types interpret as
dependent sums (left adjoint to pullback), $\Pi$-types as dependent products
(right adjoint to pullback), and identity types as the relevant path objects;
functoriality of $\llbracket-\rrbracket$ is precisely the statement that
substitution is interpreted by pullback and composes, i.e.\ soundness of the
substitution calculus. Since functors compose (\cref{thm:pipeline-functor}),
$R\circ\llbracket-\rrbracket$ is a functor. The three readings are three names
for the values of this one functor: the logical reading names the objects of
$\mathrm{Syn}(T)$ as propositions, the computational reading names its morphisms
as programs, and the physical reading names $R$ applied to them as realizations.
The equalities among these readings are literal equalities of the functor's
values, which is the content of the ``theorem $=$ specification, proof $=$
implementation, physical process $=$ realization'' slogan.
\end{proof}

\begin{corollary}[Composition is quadruply the same law]
\label{cor:composition}
Under \cref{thm:chl}, given $p:A\to B$ and $q:B\to C$, the composite $q\circ p$ is
simultaneously: logical-implication composition, proof/term composition,
program composition, and---after $R$---composition of physical processes. The
associativity \eqref{eq:assoc} of $\circ$ is one law shared by all four readings.
\end{corollary}

\begin{proof}
Immediate from \cref{thm:chl}: composition in $\mathrm{Syn}(T)$ is substitution
composition, sent by the two functors to composition in $\mathcal{C}$ and then in
$\Phys$; each named reading is the same morphism under a different name, and
associativity is inherited from \cref{def:category}.
\end{proof}

\subsection{T3: no cloning forces encoding}

\begin{theorem}[Structural no-cloning]
\label{thm:nocloning}
Let $\mathcal{C}$ be a symmetric monoidal category and suppose there is a
\emph{uniform} copying operation: writing $\mathrm{D}:\mathcal{C}\to
\mathcal{C}\times\mathcal{C}$ for the diagonal functor, a natural transformation
$\Delta:\mathrm{Id}_{\mathcal{C}}\Rightarrow\otimes\circ\mathrm{D}$ with
components $\Delta_A:A\to A\otimes A$, together with a natural deleting
$\diamond_A:A\to I$, making each $A$ a cocommutative comonoid compatibly with
$\otimes$ and the symmetry. Then $\otimes$ is a categorical product: $\mathcal{C}$
is cartesian monoidal. Consequently, in $\FdHilb$---where the tensor product is
\emph{not} a categorical product---no such uniform copying exists, and quantum
high availability cannot be replication.
\end{theorem}

\begin{proof}
The hypotheses are exactly Fox's theorem: a symmetric monoidal category equipped
with a natural, coherent, cocommutative comonoid structure
$(\Delta_A,\diamond_A)$ on every object, compatible with $\otimes$, is cartesian,
with $\otimes$ the product, $I$ the terminal object, $\diamond_A$ the unique map
$A\to I$, and $\Delta_A$ the diagonal; naturality of $\Delta$ and $\diamond$
supplies the universal property of the product on morphisms. Now suppose, for
contradiction, that $\FdHilb$ admitted such a uniform copying. Then its monoidal
product $\otimes$ would be the categorical product. But the categorical product
(equivalently, in $\FdHilb$, the biproduct) of Hilbert spaces is the direct sum
$\oplus$, with $\dim(A\oplus B)=\dim A+\dim B$, whereas
$\dim(A\otimes B)=\dim A\cdot\dim B$. For $\dim A=\dim B=2$ these are $4$ and $4$
only coincidentally; taking $\dim A=\dim B=3$ gives $6\ne 9$, so
$\otimes\not\cong\times$ as bifunctors. This contradiction shows no uniform
natural copying exists in $\FdHilb$: the no-cloning theorem. Because copying is
unavailable, protecting a quantum state cannot proceed by replication; the only
representational option is an embedding
$\iota:\mathcal{H}_{\mathrm{logical}}\hookrightarrow\mathcal{H}_{\mathrm{phys}}$
into a larger physical Hilbert space---encoding---which is precisely Part~VI's
quantum high-availability construction
$\mathcal{H}_{\mathrm{phys}}\simeq
\mathcal{H}_{\mathrm{logical}}\otimes\mathcal{H}_{\mathrm{gauge}}
\oplus\mathcal{H}_{\mathrm{error}}$.
\end{proof}

\begin{remark}[Classical structures as the copiable objects]
\label{rem:classical}
The objects of a dagger-compact category that \emph{do} admit natural copying are
exactly the \emph{classical structures}: special commutative dagger-Frobenius
algebras, whose comultiplication is a copying map. In $\FdHilb$ these correspond
to a choice of orthonormal basis---the ``classical'' pointer states one may copy.
Thus T3 also explains \emph{why} classical high availability (replication) works
for classical data and fails for genuinely quantum data: only basis states are
copiable \cite{AbramskyCoecke2004}.
\end{remark}

\subsection{T4: groupoid quotients are finer than set quotients}

\begin{theorem}[Truncation reconstruction of ``stacks retain gauge information'']
\label{thm:truncation}
Let $G$ be a $0$-type (a set carrying a group structure) acting on a $0$-type
(set) $X$. Form the \emph{groupoid quotient} (a higher inductive type)
$X/\!/G$, whose points are those of $X$, with a path $x=g\cdot x$ for each
$g:G$, subject to the groupoid coherences. Then:
\begin{enumerate}
\item[(a)] $X/\!/G$ is a $1$-type, and for each $x$ the loop space
$\Omega_x(X/\!/G)\simeq\Stab_G(x)$ retains the stabilizer as nontrivial
automorphisms;
\item[(b)] the set quotient $X/G$ is the $0$-truncation
$\lVert X/\!/G\rVert_0$, which forgets exactly those stabilizer loops;
\item[(c)] hence whenever some stabilizer is nontrivial, $X/\!/G$ is strictly
finer than $X/G$: there is no equivalence $X/\!/G\simeq X/G$.
\end{enumerate}
This is a purely type-theoretic reconstruction of the source proposition
``stacks retain gauge information'' ($\Aut_{[X/G]}(x)\simeq\Stab_G(x)$) and of
Part~III's finer-moduli phenomenon.
\end{theorem}

\begin{proof}
(a) By construction the only generating paths of $X/\!/G$ at a point $x$ are the
$g\cdot(-)$ loops, and the higher inductive coherences make loop concatenation
match group multiplication; the loops fixing $x$ are precisely those $g$ with
$g\cdot x=x$, i.e.\ $\Stab_G(x)$, giving
$\Omega_x(X/\!/G)\simeq\Stab_G(x)$. Since $G$ is a $0$-type, each stabilizer
$\Stab_G(x)\subseteq G$ is a set, so every loop space is a set and $X/\!/G$ is
$1$-truncated; here the set hypotheses on $X$ and $G$ are exactly what bounds
$X/\!/G$ at truncation level~$1$ (a $G$ with higher homotopy, such as a delooped
Lie group, would raise the level regardless of stabilizers). (b) The set
quotient is defined as the
$0$-truncation of $X/\!/G$ (equivalently the quotient of $X$ by the orbit
equivalence relation), and $0$-truncation identifies all parallel paths, so all
stabilizer loops collapse to reflexivity, retaining only orbits. (c) If some
$\Stab_G(x)\ne\{e\}$ then $\Omega_x(X/\!/G)$ is a nontrivial set while the
corresponding loop space in the $0$-type $X/G$ is contractible; an equivalence
$X/\!/G\simeq X/G$ would induce an equivalence on loop spaces (equivalences
preserve all homotopy structure), a contradiction. Reading a stabilizer loop as a
residual gauge automorphism of the configuration $x$, (a)--(c) are exactly the
statement $\Aut_{[X/G]}(x)\simeq\Stab_G(x)$ with the coarse quotient forgetting
$\Aut$, which is the source's ``stacks retain gauge information.''
\end{proof}

\begin{corollary}[One fact, four proofs]
\label{cor:four-proofs}
The proposition ``passing to the stack/groupoid quotient retains stabilizer data
that the coarse quotient forgets'' now has a type-theoretic proof
(\Cref{thm:truncation}) via truncation levels, complementing Part~I's informal
gluing argument, Part~III's algebro-geometric moduli argument, and Part~VI's
forthcoming descent-theoretic proof. These are one theorem with four proofs at
increasing rigor---the headline ``closure'' motif of the synthesis paper.
\end{corollary}

The truncation tower relating the two quotients is:
\begin{equation}
\label{eq:trunc-cd}
\begin{tikzcd}[column sep=large, row sep=large]
X \arrow[r, "\text{quotient}"] \arrow[dr, "\text{orbit map}"'] & X/\!/G \arrow[d, "\lVert-\rVert_0"] \\
& X/G
\end{tikzcd}
\end{equation}

% =====================================================================
\section{The category-theory and HoTT representation dictionary}
\label{sec:dictionary}

We collect the core dictionary entries with S/H/P labels; the full longtables are
in \cref{app:library}. The discipline is that of the whole series: an entry is
$\mathrm{S}$ when it is a standard mathematical fact used verbatim in
mathematical physics, and $\mathrm{S/H}$ or $\mathrm{H}$ when its \emph{physical}
reading is a heuristic layer.

\begin{center}
\renewcommand{\arraystretch}{1.2}
\begin{tabular}{@{}L{4.4cm}L{4.6cm}c@{}}
\toprule
\textbf{Category theory} & \textbf{Physical representation} & \textbf{Status}\\
\midrule
Object & Physical system / state space & S/H\\
Morphism $f:A\to B$ & Process, transition, evolution & S\\
Composition $g\circ f$ & Sequential process & S\\
Identity morphism & Do-nothing process & S\\
Isomorphism & Reversible equivalence & S\\
Equivalence of categories & Same theory, dual descriptions & S\\
Functor & Physical theory / quantization & S\\
Natural transformation & Uniform field redefinition & S/H\\
Monoidal product $\otimes$ & Parallel composition of systems & S\\
Dagger-compact category & Quantum process calculus & S\\
Trace & Feedback / partition function & S/H\\
Yoneda lemma & System known by its probes & S/H\\
\bottomrule
\end{tabular}
\end{center}

\begin{center}
\renewcommand{\arraystretch}{1.2}
\begin{tabular}{@{}L{4.4cm}L{4.6cm}c@{}}
\toprule
\textbf{HoTT / type theory} & \textbf{Physical representation} & \textbf{Status}\\
\midrule
Type $A$ & Space of states / possibilities & S/H\\
Term $a:A$ & State / event / witness & S/H\\
Dependent type $B(x)$ & Field / fibre over base point & S/H\\
Identity type $a=b$ & Path / equivalence of states & S/H\\
Transport & Parallel transport / coordinate change & S/H\\
Equivalence $A\simeq B$ & Duality / gauge-equivalent description & H\\
Univalence & Equivalent systems identified & S/H\\
$0$-truncation & Classical coarse-graining & H\\
Groupoid quotient $X/\!/G$ & Gauge configuration space & S/H\\
Circle HIT $S^1$ & Winding / phase sector & H\\
Cohesive HoTT & Synthetic geometry for fields & S/H\\
\bottomrule
\end{tabular}
\end{center}

% =====================================================================
\section{Haskell formalization}
\label{sec:haskell}

The package \texttt{src/category-theory-hott-composition/} encodes the
composition/identity structures of \cref{sec:cat,sec:monoidal} and the
groupoid-vs-set quotient of \Cref{thm:truncation}. We reproduce the type-class
skeleton; the runnable \texttt{Main.hs} exercises the identity/associativity
laws, a functor, a natural-transformation naturality check, a dagger-compact
no-cloning obstruction, and the two quotients.

\begin{quote}\footnotesize\ttfamily
class Category cat where\\
\hspace*{2em}identity :: cat a a\\
\hspace*{2em}compose\ \ :: cat b c -> cat a b -> cat a c\\[0.4em]
class (Category c, Category d) => CFunctor c d f where\\
\hspace*{2em}fmapC :: c a b -> d (f a) (f b)\\[0.4em]
newtype NatTrans d f g = NatTrans \char123{} component :: forall a. d (f a) (g a) \char125
\end{quote}
Here \texttt{NatTrans} is parameterized by the \emph{target category}
\texttt{d}, so a natural transformation's components are morphisms in
\texttt{d}---not hardcoded Haskell functions---which keeps the encoding honest
for non-\texttt{Hask} targets such as $\FdHilb$ or $\Bord_n$.

\texttt{Main.hs} checks numerically that \texttt{compose identity f == f == compose
f identity} for a concrete category (functions), that a chosen functor preserves
composition, that a naturality square commutes, that the tensor-vs-biproduct
dimension mismatch of \Cref{thm:nocloning} holds ($2\cdot2=4$ but $3\cdot3=9\ne
6$), and that the groupoid quotient retains a stabilizer that the set quotient
discards. The whole package compiles with \texttt{ghc}.

% =====================================================================
\section{Lean formalization (best-effort sketch)}
\label{sec:lean}

The directory \texttt{lean/category-theory-hott-composition/} contains an
idiomatic Lean~4 sketch capturing the category, functor, and identity-law types,
plus signatures for T1 and T4. It reuses Mathlib's \texttt{CategoryTheory} where
convenient and states the physics-facing wrappers (a \texttt{PhysicalCategory}
interpretation, a univalence-respects-realization signature, and a
groupoid-quotient-finer-than-set-quotient signature). It is a best-effort sketch,
not build-gated.

% =====================================================================
\section{Discussion}
\label{sec:discussion}

\paragraph{Composition in the ladder.}
Parts~I--IV produced translations---abstraction, realization, observation, TQFT,
homology---each of which is a functor. \Cref{thm:pipeline-functor} shows the
library was always a diagram in $\Cat$; Part~V simply names the ambient grammar.
The associativity and unit laws \eqref{eq:assoc}--\eqref{eq:unit} are the reason
the pipeline may be assembled from stages at all, and \Cref{cor:composition}
shows the same law governs logical, computational, and physical composition.

\paragraph{Identity in the ladder.}
Part~I \emph{postulated} that equivalent representations carry the same physical
content (Axiom~Equivalence). \Cref{thm:univalence} \emph{derives} it for every
internally definable realization, and \Cref{rem:gauge-invariance} reads the
residue---realizations that fail it---as exactly the gauge-dependent quantities.
This is the module's central payoff: identity of structures, formalized by
univalence, is the correct home for the ``same physics, different description''
principle that recurs through the series.

\paragraph{Forecasting Part~VI.}
\Cref{thm:nocloning} forces quantum high availability to be encoding, and
\Cref{thm:truncation} produces the stabilizer-retention phenomenon inside the
truncation hierarchy. Both are the categorical/type-theoretic inputs Part~VI
turns into stacks (groupoid-valued sheaves) and quantum error-correcting codes.
\Cref{cor:four-proofs} records that the stabilizer-retention fact will then have
four proofs at increasing rigor---the synthesis paper's ``closure theorem.''

\paragraph{Limitations.}
Category theory and HoTT are status~S mathematics; their \emph{physical} readings
carry the flags of \cref{sec:dictionary}. We reiterate the series limitations in
this module's terms. (i)~We claim no single universal functor $\Math\to\Phys$;
\Cref{thm:pipeline-functor} is domain-indexed by $d$, matching Part~I's
deliberate locality. (ii)~The reading of a type as a ``space of physical
possibilities'' is heuristic (status~$\mathrm{H}$). (iii)~\Cref{thm:univalence}
requires a genuinely univalent model and an internally definable realization;
externally chosen realizations are outside its scope by design, which is the
gauge-dependence caveat, not a defect. (iv)~The no-cloning argument is stated for
$\FdHilb$; infinite-dimensional refinements need additional analytic hypotheses.
(v)~Univalence is a foundational axiom (with computational content in cubical
type theory) rather than a physical measurement, so its physical reading remains
$\mathrm{S/H}$.

% =====================================================================
\section{Conclusion}
\label{sec:conclusion}

Part~V isolates the compositional glue of the representation library. Category
theory supplies the laws of \emph{composition}: associativity, units,
functoriality, naturality, the interchange law, and (via Eckmann--Hilton) the
first constraint linking two compositions through a shared identity. HoTT
supplies the theory of \emph{identity}: identity types as paths, path induction
as the sole generator of transport, and univalence as the identification of
identity with equivalence. Four theorems give the physics. Univalence turns
Part~I's Equivalence axiom into a theorem (T1); the Curry--Howard--Lambek
correspondence exhibits one functor interpreting logic, code, and physics (T2); a
structural no-cloning theorem forces quantum high availability to be encoding
(T3); and the groupoid quotient's strict refinement of the set quotient
reconstructs ``stacks retain gauge information'' inside the truncation hierarchy
(T4). Together they make explicit the universal grammar that Parts~I--IV used
implicitly and that Part~VI will use to build stacks and quantum codes. In the
compressed slogans of the series, mathematics is the syntax of physical
representation and categories/types are the grammar of its composition and
identity.

% =====================================================================
\begin{thebibliography}{99}

\bibitem{MacLane1998}
S.~Mac Lane,
\emph{Categories for the Working Mathematician},
Graduate Texts in Mathematics~5, 2nd ed., Springer, 1998.

\bibitem{HoTT2013}
The Univalent Foundations Program,
\emph{Homotopy Type Theory: Univalent Foundations of Mathematics},
Institute for Advanced Study, 2013.

\bibitem{BaezStay2009}
J.~C.~Baez and M.~Stay,
\emph{Physics, topology, logic and computation: A Rosetta Stone},
arXiv:0903.0340 (2009); \emph{Lecture Notes in Physics} \textbf{813} (2010), 95--172.

\bibitem{AbramskyCoecke2004}
S.~Abramsky and B.~Coecke,
\emph{A categorical semantics of quantum protocols},
arXiv:quant-ph/0402130 (2004); Proc.\ 19th IEEE Symposium on Logic in Computer Science (LiCS), 2004.

\bibitem{LambekScott1986}
J.~Lambek and P.~J.~Scott,
\emph{Introduction to Higher Order Categorical Logic},
Cambridge Studies in Advanced Mathematics~7, Cambridge University Press, 1986.

\bibitem{Lurie2009}
J.~Lurie,
\emph{On the classification of topological field theories},
arXiv:0905.0465 (2009); \emph{Current Developments in Mathematics 2008}, 129--280.

\bibitem{LurieHTT2009}
J.~Lurie,
\emph{Higher Topos Theory},
Annals of Mathematics Studies~170, Princeton University Press, 2009.

\bibitem{Schreiber2013}
U.~Schreiber,
\emph{Differential cohomology in a cohesive $\infty$-topos},
arXiv:1310.7930 (2013).

\bibitem{RiehlVerity2022}
E.~Riehl and D.~Verity,
\emph{Elements of $\infty$-Category Theory},
Cambridge Studies in Advanced Mathematics, Cambridge University Press, 2022.

\bibitem{Atiyah1988}
M.~Atiyah,
\emph{Topological quantum field theories},
\emph{Publications Math\'ematiques de l'IH\'ES} \textbf{68} (1988), 175--186.

\bibitem{Riehl2017}
E.~Riehl,
\emph{Category Theory in Context},
Aurora: Dover Modern Math Originals, Dover Publications, 2017.

\end{thebibliography}

% =====================================================================
\appendix
\section{Full category-theory and HoTT representation library}
\label{app:library}

The following longtables reproduce the module's dictionary (source Appendix~E),
with epistemic status labels carried verbatim. Each row is a representation entry
$(M,P,\tau,\sigma)$ in the sense of \cref{def:entry}.

\renewcommand{\arraystretch}{1.15}
\begin{longtable}{@{}L{3.9cm}L{4.3cm}L{3.2cm}c@{}}
\caption{Category theory library (physical representation dictionary).}\label{tab:ct}\\
\toprule
\textbf{Category theory} & \textbf{Physical representation} & \textbf{Semantic meaning} & \textbf{Status}\\
\midrule
\endfirsthead
\multicolumn{4}{c}{\tablename\ \thetable\ -- continued}\\
\toprule
\textbf{Category theory} & \textbf{Physical representation} & \textbf{Semantic meaning} & \textbf{Status}\\
\midrule
\endhead
\midrule
\multicolumn{4}{r}{\emph{continued on next page}}\\
\endfoot
\bottomrule
\endlastfoot
Object & Physical system / state space & Thing capable of processes & S/H\\
Morphism $f:A\to B$ & Process, transition, evolution & Physical transformation & S\\
Composition $g\circ f$ & Sequential process & One process after another & S\\
Identity morphism & Do-nothing process & Persistence / trivial evolution & S\\
Isomorphism & Reversible equivalence & Same structure, new coordinates & S\\
Equivalence of categories & Same theory structurally & Dual physical descriptions & S\\
Functor & Translation between worlds & Physical theory / quantization & S\\
Natural transformation & Uniform deformation & Theory equivalence / field redefinition & S/H\\
Category of states & State-system universe & Objects systems, arrows evolutions & H\\
Category of observables & Measurement-algebra universe & Arrows observable transformations & H\\
Monoidal category & Parallel composition & Composite systems & S\\
Tensor product $\otimes$ & Combine systems & Parallel physical composition & S\\
Unit object $I$ & Vacuum / trivial system & No-system baseline & S/H\\
Symmetric monoidal & Exchangeable systems & Bosonic-like swap symmetry & S/H\\
Braided monoidal & Nontrivial exchange & Anyonic / braided statistics & S\\
Ribbon category & Braiding plus twist & Topological quantum computation & S\\
Compact closed & Dualizable inputs/outputs & Process--state duality & S\\
Dagger category & Reversal / adjoint & Quantum adjoint, time reversal & S\\
Dagger-compact & Quantum process calculus & Teleportation, cups/caps & S\\
String diagram & Diagrammatic process & Circuit / Feynman grammar & S\\
Trace & Feedback loop & Closed process, partition function & S/H\\
Dual object $A^{*}$ & Antisystem / dual space & Bra/ket, particle/antiparticle & S/H\\
Evaluation map & Pairing & Measurement, annihilation & S/H\\
Coevaluation map & Pair creation & Entangled pair / vacuum insertion & S/H\\
Limit & Universal compatible solution & Constraint satisfaction & S\\
Colimit & Universal gluing & Global system from parts & S\\
Pullback & Fibered constraint & Simultaneous conditions & S\\
Pushout & Gluing along boundary & Coupled systems & S\\
Adjunction $F\dashv G$ & Dual translation pair & Free/forgetful, prep/measure & S/H\\
Monad & Dynamics with effects & State update, computation effect & S/H\\
Comonad & Context extraction & Observation / coarse environment & H\\
Algebra over monad & System obeying effects & Dynamics closed under update & H\\
Coalgebra & State-transition system & Observable unfolding / branching & S/H\\
Operad & Multi-input composition law & Interaction vertices / couplings & S\\
PROP & Multi-in/out network & Circuit / Feynman process algebra & S\\
Enriched category & Structured hom-spaces & Distances, probabilities, Hilbert homs & S/H\\
2-category & Objects, processes, transformations & Theories, defects, transformations & S\\
$\infty$-category & All higher coherences & Histories, higher gauge data & S\\
Higher category & Nested process hierarchy & Particles, strings, branes, defects & S\\
$n$-category & $n$-level process system & Extended TQFT codimension & S\\
Topos & Universe of variable sets & Contextual world of observation & H\\
Sheaf topos & Local-to-global universe & Observer-dependent gluing & S/H\\
Internal language & Logic of a category & Native reasoning of a universe & H\\
Yoneda lemma & Object known by probes & System determined by observations & S/H\\
\end{longtable}

\begin{longtable}{@{}L{3.9cm}L{4.3cm}L{3.2cm}c@{}}
\caption{Homotopy type theory library (physical representation dictionary).}\label{tab:hott}\\
\toprule
\textbf{HoTT / type theory} & \textbf{Physical representation} & \textbf{Semantic meaning} & \textbf{Status}\\
\midrule
\endfirsthead
\multicolumn{4}{c}{\tablename\ \thetable\ -- continued}\\
\toprule
\textbf{HoTT / type theory} & \textbf{Physical representation} & \textbf{Semantic meaning} & \textbf{Status}\\
\midrule
\endhead
\midrule
\multicolumn{4}{r}{\emph{continued on next page}}\\
\endfoot
\bottomrule
\endlastfoot
Type $A$ & Space of states / propositions & Domain of possible inhabitants & S/H\\
Term $a:A$ & State / event / witness & Actual instance of a possibility & S/H\\
Dependent type $B(x)$ & Fibre over configuration & Field degrees over base point & S/H\\
$\Sigma$-type & Total space & Configuration plus dependent data & S\\
$\Pi$-type & Space of sections & Field assigning data everywhere & S\\
Identity type $a=b$ & Path from $a$ to $b$ & Equivalence / evolution of states & S/H\\
Higher identity type & Homotopy between paths & Coherence between evolutions & S/H\\
Path induction & Reasoning from identity & Transport laws from equivalence & S/H\\
Transport & Move structure along path & Parallel transport / coordinate change & S/H\\
Equivalence $A\simeq B$ & Same physics, new representation & Duality / gauge equivalence & H\\
Univalence & Equivalence treated as identity & Equivalent systems identified & S/H\\
Universe $\mathcal U$ & Type of systems / theories & Meta-space of descriptions & H\\
Higher inductive type & Generated points/paths/cells & Synthetic topological spaces & S/H\\
Circle HIT $S^1$ & Generated loop space & Winding sector / phase circle & H\\
Sphere HIT $S^n$ & Generated $n$-charge carrier & Higher topological charge & H\\
Pushout type & Glued system & Boundary coupling / fusion & S/H\\
Pullback type & Constraint-matched system & Simultaneous compatibility & S/H\\
Quotient type & Identified states & Gauge quotient / equivalence classes & H\\
Truncation & Forget higher identity & Classical approximation & H\\
Mere proposition & Yes/no observable & Proof-irrelevant physical fact & H\\
Set-truncation & Classical state set & No higher path information & H\\
Groupoid level & Gauge-state structure & States plus equivalences & S/H\\
$\infty$-groupoid & Full homotopy state space & States, paths, paths-between-paths & S/H\\
Contractible type & Unique state up to equivalence & Trivial phase / vacuum sector & H\\
Fiber sequence & Structured dependency & Gauge fibre over base & S/H\\
Modal type theory & Observation / filtering modes & Coarse-graining, localization & H\\
Cohesive HoTT & Synthetic smooth geometry & Internal language for fields & S/H\\
Shape modality & Underlying homotopy type & Forget geometry, keep topology & S/H\\
Flat modality & Discrete / codified form & Locally constant representation & H\\
Sharp modality & Codiscrete / crisp form & Purely formal / global representation & H\\
Differential cohesion & Smooth infinitesimal structure & Synthetic differential geometry & S/H\\
\end{longtable}

\end{document}
